\title[Universal determinant completion and boundary laws]{Universal determinant completion, effective Pieri multiplication, and intrinsic primary boundary laws in multiplication failure}
\author{Qian Qi}
\date{September 23, 2026 \quad (Revision 128)}
\subjclass[2020]{14M12, 14J28, 13C40, 14B05, 14D20, 05E05}
\keywords{Multiplication failure, Fitting scheme, determinant completion, Pieri multiplication, nilpotent filtration, primary specialization, quartic K3 surface}
\hypersetup{pdftitle={Universal determinant completion, effective Pieri multiplication, and intrinsic primary boundary laws in multiplication failure},pdfauthor={Qian Qi}}
\begin{abstract}
We study the nonreduced Fitting scheme on a Grassmannian on which
multiplication fails for cube-zero algebras with Hilbert function
\((1,4,6)\).  The first nilpotent layer reconstructs the ramification
quartic K3 surface together with its quartic polarization.  The present
revision resolves the two genericity/nonvanishing gaps in the preceding
boundary atlas and replaces the former fibrewise-associated-prime sketch
by a relative-assassin argument.

The determinant bound is a special case of a general principle.  If
\(V\) has dimension \(e\), \(r\ge1\), and
\(\gamma:\Sym^rV\twoheadrightarrow S\), then for the generic
endomorphism \(M\) and \(J=I_{\dim S}(\gamma\Sym^rM)\),
\[
 (\det M)^{\binom{e+r-1}{r-1}}\in J.
\]
The exponent is sharp in the universal class.  For \(e=4,r=2\) this
gives \(d^5\in J\), hence the intrinsic nilradical has sixth power
zero on every projection-rank chart.

At projection corank four we prove the stronger membership \(d^4\in J\)
by the actual multiplication law in the matrix coordinate ring, not by
representation occurrence alone.  Standard bideterminant straightening
shows that the Cauchy components indexed by
\((4,4,4,0)\) and \((4,0,0,0)\) multiply with nonzero projection
to \((4,4,4,4)\).  The Jacobian-quartic component is nonzero on
\(G_4^\circ\); the unique \((4,4,4,4)\)-line is generated by
\(d^4\).

At projection corank two, genericity of the nine mixed-kernel rows is
proved by stabilizer slices and Groebner bases over the corresponding
rational function fields.  The same method treats every generic
\(A_H\)-rank-drop type.  It identifies the nonzero graded supports
as Segre quadrics, reduced unions of two or four ruling lines, or
explicit vertex-primary schemes, including the embedded vertex that
appears in the rank-\((1,3)\) row.  We also give explicit one-parameter
degenerations between the mixed-kernel orbit closures and record the
resulting births of ruling and embedded components.

Finally, after first imposing flatness on the finite presentations that
compute the residual colons, relative assassins and Noetherian induction
give a finite stratification on which fibrewise associated points,
minimal-versus-embedded status, incidence, generic lengths, Hilbert
polynomials, and graded multiplication ranks are constant.  This
separates the intrinsic primary signature from the coordinate-dependent
Macaulay matrices used to compute it.  Deformation to the normal cone
is retained only as the standard packaging of this intrinsic graded
algebra.
\end{abstract}
\maketitle
